% TOP_PROBLEM: {"id": 2279, "title": "Erdős Problem #724", "category_id": 1, "category": {"id": 1, "name": "number_theory", "display_name": "Number Theory", "description": "Properties of integers, prime numbers, Diophantine equations.", "slug": "number-theory", "order_index": 1, "created_at": "2026-07-31T15:26:25.670Z"}, "status": "open"}
% FIRST_POSTED: null
% ENTRY_KIND: statement_only
% Imported from: unsolvedmath_additional_500.tex; UnsolvedMath ID 2279
% !TeX program = lualatex
% Compile twice: lualatex 2279.tex
% Self-contained: no external bibliography, images, or \input files.
% Numeric section labels and visible numbers are original UnsolvedMath IDs.
\documentclass[11pt,a4paper]{article}
\usepackage[margin=24mm,headheight=14pt]{geometry}
\usepackage{fontspec}
\setmainfont{Latin Modern Roman}
\setsansfont{Latin Modern Sans}
\setmonofont{Latin Modern Mono}
\usepackage{amsmath,amssymb,mathtools}
\usepackage{unicode-math}
\setmathfont{Latin Modern Math}
\usepackage{microtype}
\usepackage{enumitem,xurl,needspace}
\usepackage[dvipsnames]{xcolor}
\usepackage{hyperref}
\hypersetup{unicode=true,colorlinks=true,linkcolor=MidnightBlue,urlcolor=MidnightBlue,
 pdftitle={UnsolvedMath Problem 2279},pdfauthor={Source-annotated compilation},bookmarksnumbered=true}
\usepackage{bookmark,tocloft,titlesec,fancyhdr}
\setlength{\cftsecnumwidth}{4.2em}
\cftsetpnumwidth{2.5em}
\cftsetrmarg{4em}
\titleformat{\section}{\Large\bfseries\raggedright}{\thesection}{0.7em}{}
\pagestyle{fancy}\fancyhf{}
\fancyhead[L]{\small\sffamily Additional UnsolvedMath statements}
\fancyhead[R]{\small Original catalogue IDs}
\fancyfoot[C]{\thepage}
\setlength{\parindent}{0pt}
\setlength{\parskip}{5pt plus 1pt}
\setlength{\emergencystretch}{3em}
\setcounter{tocdepth}{1}\setcounter{secnumdepth}{1}
\setlist[itemize]{leftmargin=3em,itemsep=4pt,topsep=4pt,parsep=0pt}
\allowdisplaybreaks
\newcommand{\N}{\mathbb{N}}
\newcommand{\Z}{\mathbb{Z}}
\newcommand{\Q}{\mathbb{Q}}
\newcommand{\R}{\mathbb{R}}
\newcommand{\C}{\mathbb{C}}
\newcommand{\F}{\mathbb{F}}
\newcommand{\A}{\mathbb{A}}
\newcommand{\PP}{\mathbb{P}}
\newcommand{\E}{\mathbb{E}}
\newcommand{\eps}{\varepsilon}
\newcommand{\dd}{\,\mathrm d}
\newcommand{\sref}[2]{\hyperlink{source:#1:#2}{[S#2]}}
\newif\ifshowcatalogue
\showcataloguetrue
\newcommand{\cataloguescope}[1]{\ifshowcatalogue
 \paragraph{Statement recorded in the supplied source.}
 #1\fi}
\begin{document}
% UnsolvedMath ID 2279; source catalogue code EP-724

\setcounter{section}{2278}

\section{Square-Root Many Mutually Orthogonal Latin Squares}\label{2279}

{\small\sffamily UnsolvedMath ID 2279 · Number Theory}

\subsection{Definitions and mathematical statement}

\paragraph{Shared notation.}Unless a section says otherwise, $\N=\{1,2,\ldots\}$,
$\Z$ is the ring of integers, $\Q,\R,\C$ are the rational, real, and complex
fields, $[N]=\{1,\ldots,\lfloor N\rfloor\}$, and $\log$ is the natural
logarithm. For real functions of a parameter tending to infinity,
$f\ll g$ and $f=O(g)$ mean $|f|\leq Cg$ eventually for some fixed $C$;
$f\gg g$ reverses this comparison; $f=o(g)$ means $f/g\to0$;
$f\sim g$ means $f/g\to1$; and $f\asymp g$ means both order bounds.
A subscript on $O$ or $\ll$ specifies allowed constant dependence.
The expression $n^{o(1)}$ in an upper bound means growth smaller than
$n^\epsilon$ for each fixed $\epsilon>0$. Local definitions govern any
exception, finite-field convention, or use of the same symbol for another
object. An asymptotic claim is not automatically an effective algorithm.



A Latin square is an $n\times n$ array using $n$ symbols once per row and column. Two squares are orthogonal if their overlaid ordered symbol pairs are all distinct. Mutual orthogonality requires this for every pair of squares in the collection. \sref{2279}{1} \sref{2279}{2}

\paragraph{Statement recorded in the supplied source.}
Let $f(n)$ be the maximum number of mutually orthogonal Latin squares of order $n$. Is it true that $ f(n) \gg n^{1/2}? $ \sref{2279}{1}

\paragraph{Residual task reported by the archive.}

The embedded review (2026-08-17) records: Prove f(n)>>sqrt(n), or improve the general exponent. \sref{2279}{1}

{\small\emph{Transcription note.} Only unambiguous escape/formatting damage and nonmathematical serialization debris were repaired in the displayed source statement.}

\subsection{Short English statement}

Does every sufficiently large order admit a number of mutually orthogonal Latin squares bounded below by a fixed multiple of its square root?

\subsection{Sources}

{\small

\begin{itemize}

\item[\hypertarget{source:2279:1}{S1}] ULAM.AI, \emph{UnsolvedMath}, supplied \texttt{problems.json}, record 2279 (EP-724), statement and background. The embedded literature-review date is 2026-08-17; that is the archive’s review date, not a fresh certification. Dataset landing page: \url{https://huggingface.co/datasets/ulamai/UnsolvedMath}.

\item[\hypertarget{source:2279:2}{S2}] T. F. Bloom, \emph{Erdős Problems}, Problem 724, \url{https://www.erdosproblems.com/724}. Reference imported from the supplied record; the live problem page was not independently checked for this compilation.

\item[\hypertarget{source:2279:3}{S3}] Bose, R. C. and Shrikhande, S. S. and Parker, E. T., Further results on the construction of mutually orthogonal Latin squares and the falsity of Euler's conjecture. Canadian J. Math. (1960), 189-203. Bibliographic information imported from the supplied record.

\item[\hypertarget{source:2279:4}{S4}] Beth, Thomas, Eine Bemerkung zur Absch\"{a}tzung der Anzahl orthogonaler lateinischer Quadrate mittels Siebverfahren. Abh. Math. Sem. Univ. Hamburg (1983), 284-288. Bibliographic information imported from the supplied record.

\item[\hypertarget{source:2279:5}{S5}] Chowla, S. and Erd\H{o}s, P. and Straus, E. G., On the maximal number of pairwise orthogonal Latin squares of a given order. Canadian J. Math. (1960), 204-208. Bibliographic information imported from the supplied record.

\item[\hypertarget{source:2279:6}{S6}] Wilson, Richard M., Concerning the number of mutually orthogonal Latin squares. Discrete Math. (1974), 181-198. Bibliographic information imported from the supplied record.

\end{itemize}

}

\label{end:2279}
\end{document}
